\documentclass{article}
\usepackage[UTF8]{ctex}
\usepackage{amsmath,amssymb,amsfonts}
\usepackage{enumitem}
\usepackage{geometry}
\geometry{a4paper, margin=1in}
% === 错题本排版优化（不参与题目标识） ===
\setlength{\parindent}{0pt}
\setlength{\parskip}{0.35em plus 0.10em minus 0.08em}
\setlength{\emergencystretch}{3em}
\setlist{itemsep=0.25em, topsep=0.40em, parsep=0pt}
\allowdisplaybreaks[2]
% === 错题本排版优化结束 ===
\begin{document}
\section*{兰州大学2025年高等代数真题}
\section*{186．兰州大学2025年高等代数试卷（回忆版）}
一．解答如下问题。

\begin{itemize}

\subsection{题目1}
  \item[1．] 已知 \(x^{n-1}+x^{n-2}+\cdots+x+1 \mid f_{1}\left(x^{n}\right)+x f_{2}\left(x^{n}\right)+\cdots+x^{n-1} f_{n}\left(x^{n}\right)\) 。证明：存在常数 \(c\) 使得
\[
(x-1)^{n} \mid \prod_{i=1}^{n}\left(f_{i}(x)-c\right) .
\]

\subsection{题目2}
  \item[2．] \(f(x)=(x-1)(x-2) \cdots(x-2 n+1)+1\) ，证明：\(f(x)\) 在 \(\mathbb{Q}\) 上不可约．
\end{itemize}

二．解答如下问题。\\

\subsection{题目3}
1．求 \(\left|\begin{array}{ccccc}1 & 2 & 3 & \cdots & n \\ x & 1 & 2 & \cdots & n-1 \\ x & x & 1 & \cdots & n-2 \\ \vdots & \vdots & \vdots & & \vdots \\ x & x & x & \cdots & 1\end{array}\right|\) ．

\subsection{题目4}
2．求 \(\left|\begin{array}{ccccc}\cos x & 1 & 0 & \cdots & 0 \\ 1 & 2 \cos x & 1 & \cdots & 0 \\ 0 & 1 & 2 \cos x & \cdots & 0 \\ \vdots & \vdots & \vdots & & \\ 0 & 0 & 0 & \cdots & 2 \cos x\end{array}\right|\)

三。设 \(A\) 为 \(n\) 阶矩阵，\(\xi_{1}, \xi_{2}, \cdots, \xi_{t}\) 为 \(A X=0\) 的一个基础解系，\(\eta\) 为 \(A X=\beta(\beta \neq 0)\) 的一个解。证明：

\begin{itemize}

  \item[1．] \(\eta, \xi_{1}, \cdots, \xi_{t}\) 线性无关。

  \item[2．] \(\eta, \eta+\xi_{1}, \cdots, \eta+\xi_{t}\) 构成 \(A X=\beta\) 的线性无关的解。

  \item[3．] \(A X=\beta\) 的任意解可表示成 \(k_{0} \eta+k_{1}\left(\eta+\xi_{1}\right)+\cdots+k_{t}\left(\eta+\xi_{t}\right)\) ，其中 \(k_{0}+k_{1}+\cdots+k_{t}=1\) 。
\end{itemize}

四．解答如下问题：

\begin{itemize}

\subsection{题目8}
  \item[1．] 设 \(A\) 为 \(m \times n\) 实矩阵，\(\beta\) 为 \(m \times 1\) 实矩阵。证明：\(r\left(A^{\mathrm{T}} A\right)=r(A)\) ，及 \(A^{\mathrm{T}} A X=A^{\mathrm{T}} \beta\) 有解。
\end{itemize}

\subsection{题目9}
2．\(A\) 为半正定矩阵，\(A B=B A\) ，证明：\(A B=B A=O\) ．

五．\(A, B\) 为 \(n\) 阶实对称矩阵，求证：存在正交阵 \(P\) ，使得 \(P^{-1} A P\) 与 \(P^{-1} B P\) 为对角阵。

六。解答如下问题：

\begin{itemize}

\subsection{题目10}
  \item[1．] 设 \(V\) 上的线性变换 \(\sigma_{1}, \sigma_{2}, \sigma_{3}, \sigma_{4}\) 满足 \(\sigma_{1} \sigma_{3}+\sigma_{2} \sigma_{4}=I_{V}, \sigma_{i} \sigma_{j}=\sigma_{j} \sigma_{i}, i, j=1,2,3,4\) ．证明：
\[
\left(\sigma_{1} \sigma_{3}\right)^{-1}(0)=\sigma_{1}^{-1}(0) \oplus \sigma_{3}^{-1}(0)
\]

\subsection{题目11}
  \item[2．] \(n\) 阶方阵 \(A, B, C, D\) 满足 \(A C+B D=I_{n}\) ，且 \(A, B, C, D\) 两两可交换．证明：\(r(A B)=r(A)+\) \(r(B)-n\) ．
\end{itemize}

\begin{itemize}
  \item[七．] \(C(A)=\{X \mid A X=X A\}, A=\left(\begin{array}{lll}1 & 1 & 0 \\ 0 & 1 & 0 \\ 0 & 0 & 1\end{array}\right)\) ．

  \item[1．] 求 \(C(A)\) 的维数。

  \item[2．] \(\sigma(X)=A X\) ，求 \(\sigma^{n}\) 的表示矩阵的若尔当标准形．
  \item[八．] \(f(x, y, z)=x^{2}+a y^{2}+z^{2}+2 b x y+2 y z+2 x z=4\) 通过正交变换
\[
\left(\begin{array}{l}
x \\
y \\
z
\end{array}\right)=P\left(\begin{array}{l}
x_{1} \\
y_{1} \\
z_{1}
\end{array}\right)
\]
得到 \(f=y_{1}^{2}+4 z_{1}^{2}=4\) ．

  \item[1．] 求 \(a, b\) 的值。

  \item[2．] 求正交矩阵 \(P\) 。
\end{itemize}
\end{document}
